\documentclass{amsart}
% For dealing with references we use the comment environment
\usepackage{verbatim}
\newenvironment{reference}{\comment}{\endcomment}
%\newenvironment{reference}{}{}
\newenvironment{slogan}{\comment}{\endcomment}
\newenvironment{history}{\comment}{\endcomment}

% For commutative diagrams we use Xy-pic
\usepackage[all]{xy}

% We use 2cell for 2-commutative diagrams.
\xyoption{2cell}
\UseAllTwocells

% We use multicol for the list of chapters between chapters
\usepackage{multicol}

% This is generally recommended for better output
\usepackage{lmodern}
\usepackage[T1]{fontenc}

% For cross-file-references

% Package for hypertext links:
\usepackage{hyperref}

% For any local file, say "hello.tex" you want to link to please

% Theorem environments.
%
\theoremstyle{plain}
\newtheorem{theorem}[subsection]{Theorem}
\newtheorem{proposition}[subsection]{Proposition}
\newtheorem{lemma}[subsection]{Lemma}

\theoremstyle{definition}
\newtheorem{definition}[subsection]{Definition}
\newtheorem{example}[subsection]{Example}
\newtheorem{exercise}[subsection]{Exercise}
\newtheorem{situation}[subsection]{Situation}

\theoremstyle{remark}
\newtheorem{remark}[subsection]{Remark}
\newtheorem{remarks}[subsection]{Remarks}

\numberwithin{equation}{subsection}

% Macros
%
\def\lim{\mathop{\mathrm{lim}}\nolimits}
\def\colim{\mathop{\mathrm{colim}}\nolimits}
\def\Spec{\mathop{\mathrm{Spec}}}
\def\Hom{\mathop{\mathrm{Hom}}\nolimits}
\def\Ext{\mathop{\mathrm{Ext}}\nolimits}
\def\SheafHom{\mathop{\mathcal{H}\!\mathit{om}}\nolimits}
\def\SheafExt{\mathop{\mathcal{E}\!\mathit{xt}}\nolimits}
\def\Sch{\mathit{Sch}}
\def\Mor{\mathop{\mathrm{Mor}}\nolimits}
\def\Ob{\mathop{\mathrm{Ob}}\nolimits}
\def\Sh{\mathop{\mathit{Sh}}\nolimits}
\def\NL{\mathop{N\!L}\nolimits}
\def\CH{\mathop{\mathrm{CH}}\nolimits}
\def\proetale{{pro\text{-}\acute{e}tale}}
\def\etale{{\acute{e}tale}}
\def\QCoh{\mathit{QCoh}}
\def\Ker{\mathop{\mathrm{Ker}}}
\def\Im{\mathop{\mathrm{Im}}}
\def\Coker{\mathop{\mathrm{Coker}}}
\def\Coim{\mathop{\mathrm{Coim}}}

% Boxtimes
%
\DeclareMathSymbol{\boxtimes}{\mathbin}{AMSa}{"02}

%
% Macros for moduli stacks/spaces
%
\def\QCohstack{\mathcal{QC}\!\mathit{oh}}
\def\Cohstack{\mathcal{C}\!\mathit{oh}}
\def\Spacesstack{\mathcal{S}\!\mathit{paces}}
\def\Quotfunctor{\mathrm{Quot}}
\def\Hilbfunctor{\mathrm{Hilb}}
\def\Curvesstack{\mathcal{C}\!\mathit{urves}}
\def\Polarizedstack{\mathcal{P}\!\mathit{olarized}}
\def\Complexesstack{\mathcal{C}\!\mathit{omplexes}}
% \Pic is the operator that assigns to X its picard group, usage \Pic(X)
% \Picardstack_{X/B} denotes the Picard stack of X over B
% \Picardfunctor_{X/B} denotes the Picard functor of X over B
\def\Pic{\mathop{\mathrm{Pic}}\nolimits}
\def\Picardstack{\mathcal{P}\!\mathit{ic}}
\def\Picardfunctor{\mathrm{Pic}}
\def\Deformationcategory{\mathcal{D}\!\mathit{ef}}

\setlength{\emergencystretch}{3em}
\begin{document}
\title[Stacks Project, Tag 00R4]{112 --- Miracle flatness}
\maketitle
\noindent Copyright (C) 2005--2025 Johan de Jong. Stacks Project authors. Source: \href{https://stacks.math.columbia.edu/tag/00R4}{Tag 00R4}.

\noindent Editorial difficulty note (not author text). Difficulty H2: The dimension equality permits the choice of an element that is a regular parameter of the source and a nonzerodivisor of the Cohen-Macaulay target. Quotienting lowers both dimensions, allowing induction, and the local flatness criterion then lifts flatness from the quotient..

\noindent Editorial provenance note (not author text). History: excerpt from revision \texttt{a0444\allowbreak 6e57e\allowbreak c1fbc\allowbreak 252a8\allowbreak 71afc\allowbreak ec775\allowbreak 2fb28\allowbreak 07b14}, algebra.tex, lines 32653--32691. Reference presentation was adapted; original-author context and core support blocks were retained or added. The mathematical wording is unchanged. Licensed under GNU FDL 1.2 or later; no invariant sections or cover texts. See ../licenses/COPYING. Context blocks reproduce author definitions and supporting results; they are not additional counted problems.

\begin{lemma}
\label{lemma-CM-over-regular-flat}
\begin{slogan}
Miracle flatness
\end{slogan}
Let $R \to S$ be a local homomorphism of Noetherian local
rings. Assume
\begin{enumerate}
\item $R$ is regular,
\item $S$ Cohen-Macaulay,
\item $\dim(S) = \dim(R) + \dim(S/\mathfrak m_R S)$.
\end{enumerate}
Then $R \to S$ is flat.
\end{lemma}

\begin{proof}
By induction on $\dim(R)$. The case $\dim(R) = 0$ is trivial, because
then $R$ is a field. Assume $\dim(R) > 0$. By (3) this implies that
$\dim(S) > 0$. Let $\mathfrak q_1, \ldots, \mathfrak q_r$ be the minimal
primes of $S$. Note that $\mathfrak q_i \not \supset \mathfrak m_R S$ since
$$
\dim(S/\mathfrak q_i) = \dim(S) > \dim(S/\mathfrak m_R S)
$$
the first equality by Lemma \href{https://stacks.math.columbia.edu/tag/00N9}{lemma maximal chain CM (Tag 00N9)} and the
inequality by (3). Thus
$\mathfrak p_i = R \cap \mathfrak q_i$ is not equal to $\mathfrak m_R$.
Pick $x \in \mathfrak m_R$, $x \not \in \mathfrak m_R^2$, and
$x \not \in \mathfrak p_i$, see
Lemma \href{https://stacks.math.columbia.edu/tag/00DS}{lemma silly (Tag 00DS)}.
Hence we see that $x$ is not contained in any of the minimal
primes of $S$. Hence $x$ is a nonzerodivisor on $S$ by (2), see
Lemma \href{https://stacks.math.columbia.edu/tag/02JN}{lemma reformulate CM (Tag 02JN)} and
$S/xS$ is Cohen-Macaulay with $\dim(S/xS) = \dim(S) - 1$.
By (1) and Lemma \href{https://stacks.math.columbia.edu/tag/00NQ}{lemma regular ring CM (Tag 00NQ)} the ring $R/xR$ is regular
with $\dim(R/xR) = \dim(R) - 1$.
By induction we see that $R/xR \to S/xS$ is flat. Hence we
conclude by Lemma \href{https://stacks.math.columbia.edu/tag/00ML}{lemma variant local criterion flatness (Tag 00ML)}
and the remark following it.
\end{proof}
\end{document}
